\section{ИСТОРИЧЕСКИЕ ЭТАПЫ РАЗВИТИЯ СРАВНИТЕЛЬНОЙ ПОЛИТОЛОГИИ}

\subsection{Становление дисциплины и традиционный период}
Становление сравнительной политологии как самостоятельной отрасли знания происходило во второй половине XIX века и было тесно связано с развитием историко-сравнительного и формально-легального подходов~\cite{smorgunov2012}. В этот период значительный вклад внесли такие исследователи, как Э. Фримен, В. Вильсон и Дж. Берджес. В российском научном пространстве основы методологии сопоставления закладывались в трудах М. М. Ковалевского, который рассматривал компаративистику как посредника между теорией и эмпирикой~\cite{smorgunov2012}.

Первая половина XX века ознаменовалась господством «традиционного подхода», для которого, по классической оценке Р. Макридиса, были характерны дескриптивность, формальный легализм, методологическое безразличие и евроцентристский прогрессизм~\cite{smorgunov2012}. Сравнительные исследования этого времени преимущественно сводились к конфигуративному описанию институтов узкого круга ведущих западных держав.

\subsection{От «новой» к плюралистической и неоинституциональной компаративистике}
Рубежным этапом стал Эванстонский семинар 1952 года, заложивший основы «новой» сравнительной политологии~\cite{smorgunov2012}. Под эгидой Комитета по сравнительной политологии (SSRC) во главе с Г. Алмондом произошла поведенческая (бихевиоральная) и системная революция~\cite{smorgunov2012}. В фокус исследователей попали не только формальные нормы, но и реальное политическое поведение, а также незападные политические системы. Ключевыми достижениями этого периода стали системные модели Г. Алмонда и Б. Пауэлла, а также «концептуальная карта Европы» С. Роккана~\cite{smorgunov2012}.

С середины 1970-х годов дисциплина вступает в фазу «плюрализма», вызванную критикой бихевиорального и структурно-функционального подходов за их ценностную нейтральность и утрату уникальности контекста~\cite{smorgunov2012}. Происходит радикализация под под влиянием постмодернистских и феминистских волн, возрождается интерес к сравнительно-исторической методологии М. Вебера и К. Маркса~\cite{smorgunov2012}.

На рубеже XX--XXI веков ключевой парадигмой становится неоинституционализм в различных его версиях (рационального выбора, исторический, социологический)~\cite{smorgunov2012}. Неоинституциональный этап позволил усилить объяснительную функцию сравнения, сфокусировавшись на анализе условий возникновения, функционирования и трансформации политических институтов~\cite{smorgunov2012}.